La solución original calcula, para cada órbita $i$ en orden creciente, dos cosas por posición $j$: $B(i,j)$ (el mejor costo con el que se \emph{entra} a $(i,j)$ desde fuera de la órbita, ya sea por ser la entrada del campo o por un portal) y $E(i,j)$ (el mejor costo tras moverse dentro de la órbita $i$ hasta $j$). Para recuperar la ruta basta con recordar, además del costo, \emph{de dónde} vino cada uno de esos dos valores:

\begin{itemize}
  \item \texttt{origen\_externo[i][j]}: cómo se llegó a $B(i,j)$ --- \texttt{\textquotedbl ENTRADA\textquotedbl} si $(i,j) = (1,1)$, o la posición $(x_s, y_s)$ del portal que lo produjo.
  \item \texttt{origen\_interno[i][j]}: la posición $k$ de la órbita $i$ desde la que se caminó (en línea recta, en uno de los dos barridos) hasta llegar a $j$ con costo $E(i,j)$.
\end{itemize}

Ambos se llenan con el mismo trabajo $\Theta(1)$ por posición que ya hace la solución original al actualizar el mínimo acumulado de cada barrido, así que guardarlos no cambia la complejidad: sigue siendo $\Theta(nm)$ en tiempo y $\Theta(nm)$ en espacio.

Para reconstruir la ruta se parte de $(n,m)$ y se deshacen los pasos: se camina en línea recta de \texttt{origen\_interno[i][j]} a $j$ dentro de la órbita $i$, agregando cada posición intermedia, y desde ahí se salta a \texttt{origen\_externo} --- que es o bien la entrada (fin de la reconstrucción) o la posición de un portal en una órbita anterior, desde la que se repite el proceso.

\begin{pseudocode}
def energia_minima_con_ruta(n, m, e, portales):
    INF = float("inf")
    A = [[INF] * m for _ in range(n)]
    A[0][0] = 0

    origen_externo = [[None] * m for _ in range(n)]
    origen_externo[0][0] = "ENTRADA"
    origen_interno = [[None] * m for _ in range(n)]

    salientes = [[] for _ in range(n)]
    for xs, ys, xe, ye in portales:
        salientes[xs].append((ys, xe, ye))

    for i in range(n):
        fila_a, ei = A[i], e[i]
        fila_dp = [INF] * m
        interno = origen_interno[i]

        mejor, mejor_k = INF, None
        for j in range(m):
            candidato = fila_a[j] - j * ei
            if candidato < mejor:
                mejor, mejor_k = candidato, j
            fila_dp[j] = mejor + j * ei
            interno[j] = mejor_k

        mejor, mejor_k = INF, None
        for j in range(m - 1, -1, -1):
            candidato = fila_a[j] + j * ei
            if candidato < mejor:
                mejor, mejor_k = candidato, j
            candidato = mejor - j * ei
            if candidato < fila_dp[j]:
                fila_dp[j], interno[j] = candidato, mejor_k

        for ys, xe, ye in salientes[i]:
            if fila_dp[ys] < A[xe][ye]:
                A[xe][ye] = fila_dp[ys]
                origen_externo[xe][ye] = (i, ys)

    if fila_dp[m - 1] == INF:
        return INF, None

    ruta = []
    i, j = n - 1, m - 1
    while True:
        k = origen_interno[i][j]
        paso = 1 if j >= k else -1
        for pos in range(j, k - paso, -paso):
            ruta.append((i, pos))
        destino_previo = origen_externo[i][k]
        if destino_previo == "ENTRADA":
            break
        i, j = destino_previo

    ruta.reverse()
    return fila_dp[m - 1], ruta
\end{pseudocode}
